% ============================================================
% PAGE 5 — Module 1: Advanced Feature Extraction Backbone
% ============================================================

\subsection{Module 1: Advanced Feature Extraction Backbone}

The backbone is the computational engine that transforms raw
pixel intensities into hierarchically organized feature
representations. In FusionDet the CNN branch of the dual
backbone carries the bulk of the spatial feature extraction
workload, and its design therefore exerts a disproportionate
influence on both accuracy and inference latency. We adopt a
customized \textit{Cross-Stage Partial Darknet}
(CSP-Darknet) \cite{wang2020cspnet} as the convolutional
backbone, augmented with a lightweight spatial attention
mechanism that sharpens the network's focus on
informationally dense regions of each feature map.

\subsubsection{CSP-Darknet Formulation}

The defining characteristic of CSP-Darknet is its partitioning
of the input feature map at each stage into two parallel
streams: a \textit{dense stream} that passes through a
sequence of residual bottleneck blocks, and a \textit{shortcut
stream} that bypasses those blocks entirely. Let
$\mathbf{X}^{l} \in \mathbb{R}^{C \times H \times W}$ denote
the input to stage $l$. The channel dimension is split evenly:
\begin{equation}
    \mathbf{X}_{\text{dense}}^{l},\;
    \mathbf{X}_{\text{short}}^{l}
    = \text{Split}\!\bigl(\mathbf{X}^{l}\bigr)
    \label{eq:csp_split}
\end{equation}
The dense stream is processed by $M$ sequential bottleneck
blocks:
\begin{equation}
    \mathbf{Y}_{\text{dense}}^{l}
    = \bigl(\mathcal{B}_{M}^{l} \circ \cdots \circ
      \mathcal{B}_{2}^{l} \circ
      \mathcal{B}_{1}^{l}\bigr)
      \bigl(\mathbf{X}_{\text{dense}}^{l}\bigr)
    \label{eq:csp_dense}
\end{equation}
where each $\mathcal{B}_{m}^{l}$ is a residual block
comprising a $1 \times 1$ channel-reduction convolution followed
by a $3 \times 3$ depthwise-separable convolution
\cite{howard2017mobilenets} and a residual addition. The two
streams are then re-united via concatenation and a $1 \times 1$
transition convolution $\mathcal{T}^{l}$:
\begin{equation}
    \mathbf{X}^{l+1}
    = \mathcal{T}^{l}\!\Bigl(
        \text{Concat}\bigl(
            \mathbf{Y}_{\text{dense}}^{l},\;
            \mathbf{X}_{\text{short}}^{l}
        \bigr)
      \Bigr)
    \label{eq:csp_merge}
\end{equation}
This cross-stage partial connectivity yields two concrete
benefits. First, because only half the channels traverse the
expensive bottleneck chain, the floating-point operation count
is reduced by nearly 40\% relative to a plain residual stack of
equivalent depth \cite{wang2020cspnet}. Second, the gradient
signal flowing through the shortcut stream remains undiluted by
the depth of the dense path, improving convergence stability
during training.

Our implementation instantiates four CSP stages with channel
widths $\{64, 128, 256, 512\}$ and bottleneck counts
$M \in \{2, 4, 6, 4\}$, producing multi-scale feature maps
$\{P_2, P_3, P_4, P_5\}$ at strides $\{4, 8, 16, 32\}$
relative to the input resolution.

\subsubsection{Activation Functions}

The choice of activation function within each bottleneck
materially affects both representational capacity and hardware
throughput. We employ the \textit{Sigmoid Linear Unit} (SiLU),
also known as Swish-1 \cite{ramachandran2018searching}:
\begin{equation}
    \text{SiLU}(x) = x \cdot \sigma(x)
                    = \frac{x}{1 + e^{-x}}
    \label{eq:silu}
\end{equation}
where $\sigma(\cdot)$ denotes the logistic sigmoid. Unlike
ReLU, SiLU is smooth and non-monotonic---it permits small
negative activations, which has been shown to improve gradient
flow in deep residual topologies \cite{ramachandran2018searching}.
Compared with the closely related Mish activation
\cite{misra2020mish},
\begin{equation}
    \text{Mish}(x) = x \cdot \tanh\!\bigl(
        \ln(1 + e^{x})\bigr)
    \label{eq:mish}
\end{equation}
SiLU offers a marginally lower computational cost because it
avoids the nested $\tanh$--softplus composition, while
empirical benchmarks on COCO indicate that the two activations
produce statistically indistinguishable mAP when all other
hyperparameters are held constant
\cite{bochkovskiy2020yolov4}. We therefore select SiLU as the
default activation throughout the backbone, neck, and detection
heads, prioritizing the modest latency advantage on both GPU
and edge inference runtimes.

\subsubsection{Spatial Attention Augmentation}

Standard convolutional stages treat every spatial location
within a feature map with equal computational weight, yet not
all locations carry equal information---object-bearing regions
are typically far more salient than background. To bias the
backbone toward these high-value regions without incurring a
significant parameter overhead, we insert a lightweight
\textit{Spatial Attention Module} (SAM) after the final
bottleneck of each CSP stage.

Given a feature map
$\mathbf{F} \in \mathbb{R}^{C \times H \times W}$, SAM first
compresses the channel dimension by computing both a
channel-wise max-pool and a channel-wise average-pool,
yielding two single-channel descriptors:
\begin{equation}
    \mathbf{D}_{\max} = \max_{c}\, \mathbf{F}_{c},
    \qquad
    \mathbf{D}_{\text{avg}} = \frac{1}{C}
    \sum_{c=1}^{C} \mathbf{F}_{c}
    \label{eq:sam_pool}
\end{equation}
These descriptors are concatenated along the channel axis and
passed through a $7 \times 7$ convolution followed by a sigmoid
gate to produce a spatial attention map
$\mathbf{M}_{s} \in [0,1]^{H \times W}$:
\begin{equation}
    \mathbf{M}_{s} = \sigma\!\Bigl(
        f^{7 \times 7}\bigl(
            \text{Concat}(\mathbf{D}_{\max},\,
                           \mathbf{D}_{\text{avg}})
        \bigr)
    \Bigr)
    \label{eq:sam_gate}
\end{equation}
The refined feature map is then obtained via element-wise
multiplication:
\begin{equation}
    \hat{\mathbf{F}} = \mathbf{M}_{s} \odot \mathbf{F}
    \label{eq:sam_refine}
\end{equation}
Because the entire SAM pathway involves only a single
$7 \times 7$ convolution operating on a two-channel input, it
adds fewer than $0.01$~GFLOPs per stage---a negligible cost
that nonetheless provides a consistent $0.4$--$0.6$~point mAP
improvement across all scale levels in our ablation
experiments (detailed in Section~\ref{sec:results}). This
design draws on the spatial branch of CBAM
\cite{woo2018cbam} but omits the channel-attention component,
which is already handled globally by the AFF module described
in (\ref{eq:aff}).
